% Part: turing-machines
% Chapter: undecidability
% Section: universal-tm

\documentclass[../../../include/open-logic-section]{subfiles}

\begin{document}

\olfileid{tur}{und}{uni}
\olsection{通用图灵机}

在\olref[enu]{sec}节中，我们讨论了每台图灵机如何可以由一个有限的整数序列来描述。该序列编码了图灵机的状态、字母表、起始状态以及指令。我们还指出，所有这些描述的集合是可数的。既然这类描述的集合是可数无穷的，就意味着存在一个从~$\Nat$ 到这些描述的满射函数。运用康托尔之字形法，就可以得到这样一个满射函数。这为我们提供了一种枚举所有图灵机（描述）的方法。固定了这样一种枚举后，谈论第~$1$ 台、第~$2$ 台、……、第~$e$ 台图灵机便有了意义。这些数字被称为\emph{指标}。

\begin{defn}
如果 $M$ 是第~$e$ 台图灵机（在我们固定的枚举中），则称 $e$ 是~$M$ 的一个\emph{指标}。我们用 $M_e$ 表示第~$e$ 台图灵机。
\end{defn}

一台机器可能有不止一个指标，例如，对~$M$ 的两种描述可能仅在列出指令的顺序上有所不同，而这些不同的描述将具有不同的指标。

重要的是，我们可以给出一种图灵机描述的枚举，使得从指标~$e$ 能够有效计算出~$M$ 的描述，且从机器~$M$ 的描述也能有效计算出它的一个指标。根据丘奇--图灵论题，存在一台图灵机，它能够根据指标~$e$ 恢复对应图灵机的描述，并将该描述作为输出写在纸带上。该描述将是若干由~$\TMstroke$ 构成的区块序列（代表描述~$M_e$ 的序列中的正整数）。

由此自然产生了一个问题：图灵机指标的哪些函数本身可由图灵机计算？图灵机指标的哪些性质可由图灵机判定？举个例子：将指标~$e$ 映射到具有该指标的图灵机所拥有的状态数的函数，可由图灵机计算。这样一台图灵机会这样做：以纸带上包含单个由~$e$ 个 $\TMstroke$ 组成的区块作为初始状态开始，它首先将~$e$ 解码为其描述。此时，描述由纸带上若干由~$\TMstroke$ 构成的区块序列来表示。由于该序列的第一个元素就是状态数，因此现在只需擦除第一个 $\TMstroke$ 区块以外的所有内容，然后停机。

一个重要的结果如下：

\begin{thm}\ollabel{thm:universal-tm} 存在一台\emph{通用图灵机}~$U$，当以输入 $\tuple{e,n}$ 启动时：
  \begin{enumerate}
    \item $U$ 停机当且仅当 $M_e$ 在输入~$n$ 上停机，并且
    \item 如果 $M_e$ 以输出 $m$ 停机，那么~$U$ 也以输出 $m$ 停机。
  \end{enumerate}
  因此，$U$ 计算了部分函数 $f\colon \Nat \times \Nat \pto \Nat$：若 $M_e$ 以输入~$n$ 启动并以输出~$m$ 停机，则 $f(e,n) = m$；否则 $f(e,n)$ 无定义。
\end{thm}

\begin{proof}
  实际构造出~$U$ 基本上是不可能的，因为它是一台极其复杂的机器。但我们可以概述它的工作方式，然后援引丘奇--图灵论题。当~$U$ 启动时，它的纸带上包含一个由 $e$ 个 $\TMstroke$ 组成的区块，后面跟着一个由 $n$ 个 $\TMstroke$ 组成的区块。它首先将指标~$e$``解码''至输入~$n$ 的右侧。这会得到一个数字列表（即由~$\TMblank$ 分隔的多个 $\TMstroke$ 区块），该列表描述了机器~$M_e$ 的指令。接着，$U$ 在 $M_e$ 描述的右侧纸带上，写下~$M_e$ 的起始状态编号以及数字~$1$。（同样，这些都以一元形式表示，即由 $\TMstroke$ 构成的区块。）然后，它将输入（由 $n$ 个 $\TMstroke$ 组成的区块）复制到右侧——但它将每个 $\TMstroke$ 替换为一个由三个 $\TMstroke$ 组成的区块（记住，符号 $\TMstroke$ 的编码是~$3$，$\TMendtape$ 的编码是~$1$，$\TMblank$ 的编码是~$2$）。在这一系列区块（在~$U$ 的纸带上由 $\TMblank$ 符号分隔）的左端，它写入一个单独的~$\TMstroke$，这是~$\TMendtape$ 的编码。

  现在，$U$ 的纸带上记录了：指标~$e$，数字~$n$，起始状态的编码数（``当前状态''），初始读写头位置~$1$ 的编号（``当前读写头位置''），以及``纸带''的初始内容（一系列代表~$M_e$ 各符号编码数的 $\TMstroke$ 区块——即``符号''——由~$\TMblank$ 分隔）。

  它现在通过执行以下操作，来模拟~$M_e$ 在输入~$n$ 上启动后的行为：
  \begin{enumerate}
    \item 找到``当前读写头位置''的编号~$k$（开始时是~$1$），
    \item 移动到``纸带''上的第~$k$ 个区块，查看那里的``符号''是什么，
    \item\ollabel{find-inst}%
    找到与当前``状态''和``符号''匹配的指令，
    \item 移回``纸带''上的第~$k$ 个区块，并将那里的``符号''替换为~$M_e$ 会写下的符号的编码数，
    \item 将读写头移动到记录当前``状态''的地方，并将那里的数字替换为新状态的编号，
    \item 移动到记录``纸带位置''的地方，擦除一个~$\TMstroke$ 或添加一个~$\TMstroke$（若指令分别要求向左或向右移动）。
    \item 重复。\footnote{我们在此忽略了一些微妙的困难。例如，当~$U$ 增加记录``当前读写头位置''的计数器时，它可能需要一些额外的空间——在这种情况下，它将不得不把整个``纸带''向右移动。}
  \end{enumerate}
  如果~$M_e$ 在输入~$n$ 上启动后永不停机，那么~$U$ 也永不停机，因此它的输出无定义。

  如果在步骤~\olref{find-inst} 中发现~$M_e$ 的描述不包含针对当前``状态''/``符号''组合的指令，那么~$M_e$ 将会停机。若发生这种情况，$U$ 会擦除``纸带''左侧的部分。对于每一个由三个~$\TMstroke$ 组成的区块（代表~$M_e$ 纸带上的一个~$\TMstroke$），它会在自己纸带的左端写入一个 $\TMstroke$，并依次擦除``纸带''。完成后，$U$ 的纸带上包含一个长度为~$m$ 的、由~$\TMstroke$ 组成的区块。
  
  如果~$U$ 在``纸带''上遇到非三个~$\TMstroke$ 构成的区块，它会立即停机。由于此时~$U$ 的纸带并不包含单一的一个由~$\TMstroke$ 构成的区块，其输出就不是一个自然数，即，在这种情况下 $f(e,n)$ 是无定义的。
\end{proof}

\end{document}